\documentclass[11pt,aps,preprint,nofootinbib,superscriptaddress,tightenlines]{revtex4-2}
\usepackage{amsmath,amssymb,bm,mathtools,graphicx}
\usepackage{xcolor}
\definecolor{darkblue}{RGB}{0,0,160}
\usepackage[colorlinks=true,linkcolor=darkblue,citecolor=darkblue,urlcolor=darkblue]{hyperref}

\def\Sect#1{Sect.~\ref{#1}}
\newcommand{\file}[1]{\texttt{#1}}
\let\oldparagraph\paragraph
\renewcommand{\paragraph}{\par\bigskip\oldparagraph}

\begin{document}
\sloppy

\title{The code: names, defaults, run records and tests}
\author{D.~Teaney}
\affiliation{Department of Physics and Astronomy, Stony Brook University,
Stony Brook, New York 11794, USA}
\date{\today}

\begin{abstract}
A plan to put the code in order before the next campaigns: one name for the stochastic step,
the superseded code set aside, every default defined once, a run record generated from the
settings actually used, and a test suite in one place with pass/fail criteria.  None of it
changes the numerics; every step is checked bit for bit against the code before it.
\end{abstract}

\maketitle

%=============================================================================
\section{The rule: no change to the numerics}
\label{sec:rule}

Each step below is one commit on the branch \file{cleanup}, and after each:
\begin{itemize}
\item \file{test/testregress\_step.jl} (the full coupled step, thermal start, $\lambda=1$ and
  $\lambda=\infty$, twelve steps) gives $x$, $m$, $s$, the acceptances and the guard counts
  bit-identical to \file{main}, for the n-scheme (the default) and for the smoothed
  Hamiltonian, at 1 and 4 threads;
\item every unit test gives the same output as on \file{main}, timing aside;
\item a step that touches a driver also compares a short fixed-seed run of that driver bit for
  bit (its \file{.h5}).
\end{itemize}
Two steps are done (A1, A2): four functions nothing called were removed, and the PyPlot
functions moved from \file{src/setup.jl} to \file{examples/plotsetup.jl}.

%=============================================================================
\section{Names and layout}
\label{sec:names}

\paragraph{The stochastic step.}  ``stochrep'' (stochastic step in the particle
representation) becomes ``stochastic'':
\begin{center}
\begin{tabular}{ll}
\hline
now & after \\ \hline
\file{src/stochrep.jl}, \file{stochrep\_step!} & \file{src/stochastic.jl}, \file{stochastic\_step!} \\
\file{examples/stochrep\_equilibrium.jl} & \file{examples/equilibrium.jl} \\
\file{examples/stochrep\_hamiltonian.jl} & \file{examples/stability.jl} \\
other \file{examples/stochrep\_*.jl} & \file{examples/stochastic\_*.jl} \\
\file{docs/stochrep.tex} & \file{docs/stochastic.tex} \\
\hline
\end{tabular}
\end{center}
References in the notes, \file{CLAUDE.md}, \file{nersc/README.md} and \file{nersc/LOG.md}
(the current sections) are updated.  The campaign files of finished campaigns are records of
what was run and keep the old names; new campaigns use the new ones.  The rename is merged
only when no campaign is running at NERSC, since a running job array reads the driver by
name.

\paragraph{The files of the library.}  Two files are named after history rather than
content: \file{changerep.jl} (``change of representation'') holds the maps between the
particle variables $(s_a,\bm m_a)$ and the pointwise thermodynamics, plus the deterministic
viscous step built on them; \file{celllist.jl} (784 lines) holds the cell list but also the
smoothing, the ideal step, the time step and the bond dissipation.  They are split by content:
\begin{center}
\begin{tabular}{lp{0.62\textwidth}}
\hline
file & contents \\ \hline
\file{sph.jl} & kernels, \file{SimParams}, equation of state, \file{Particle} (as now) \\
\file{celllist.jl} & cell list and neighbour list only \\
\file{ideal.jl} & smoothing, velocities, Hamiltonian force, leapfrog, \file{compute\_dt} \\
\file{pointwise.jl} & particle $\leftrightarrow$ pointwise maps: \file{pointwise\_state},
  \file{particle\_energy}, \file{eos\_from\_energy}, \file{rest\_energy},
  \file{inverse\_susceptibility}, the star fit \\
\file{viscous.jl} & the deterministic viscous step: bond weights, bond dissipation,
  \file{dissipative\_step!} and its variants \\
\file{stochastic.jl} & the stochastic step (now \file{stochrep.jl}) \\
the others & \file{setup.jl}, \file{rng.jl}, \file{measure.jl}, \file{monitor.jl}, \file{runinfo.jl}, as now \\
\hline
\end{tabular}
\end{center}
The deterministic step variants lose ``changerep'' from their names:
\file{dissipative\_step\_changerep!}, \file{\_exact!}, \file{\_grad!} become
\file{viscous\_step\_first!}, \file{viscous\_step\_exact!}, \file{viscous\_step\_grad!} (first-order,
exact and gradient-corrected back maps), and \file{dissipative\_step!} becomes
\file{viscous\_step!}.  The value \file{SPH\_DISSIPATION=changerep} stays accepted (old
campaign files use it) as an alias of the new name, \file{particle}.  \file{docs/changerep.tex}
keeps its name: it is the record of that study.

\paragraph{Superseded code} moves to \file{src/legacy/}: the old particle sweep with the
noise tensor (the present \file{src/stochastic.jl}, whose name the new step takes) and the
two-level (strain) dissipation (\file{two\_level\_dissipative\_force!}, \file{compute\_strain},
\file{kappa\_dot\_strain}, \file{SPH\_DISSIPATION=two\_level}).  The bond dissipation and
\file{changerep.jl} stay: the deterministic step of the n-scheme uses the first, and the
stochastic step and the equilibrium start use the second.  The five examples
and the test that use it include it from there and still run, so the figures of the older
notes can be made again; the main library no longer loads it.

\paragraph{One load file.}  \file{src/load.jl} includes the library in order.  About seventy
scripts now repeat their own list of \file{include}s, and the lists differ.  Each becomes
\file{include} of \file{src/load.jl}.

\paragraph{The kernel.}  \file{wendland\_c2} and \file{wendland\_c2\_grad} also evaluate the
C4 and C6 kernels; they become \file{kernel\_W} and \file{kernel\_gradW} (19 files, names only).

%=============================================================================
\section{Defaults, defined once}
\label{sec:defaults}

Today a default can live in four places: a field of \file{SimParams}, a keyword default of
\file{stochrep\_step!}, a \file{get(ENV, ...)} in a driver, and the hand-written run record.
They drift apart.  After:
\begin{itemize}
\item \textbf{Structs hold the settings, with their defaults:} \file{SimParams} (scheme,
  kernel, $h$, $a$, box), the equation of state, and a new \file{StochOptions} for the
  stochastic step ($\lambda$, residual move, vote, grid shift, gradient map).  A default that
  depends on another setting (the gradient map is on except for the n-scheme) is resolved
  when the struct is built, so the struct holds the value used.
  \file{stochastic\_step!} takes a \file{StochOptions}; the keywords stay as a thin wrapper
  while the drivers and tests change over.
\item \textbf{One reader for the environment:} \file{setting("SR\_ETAS", 0.2)} returns the value
  and notes it, with its default, for the record.  Drivers stop reading \file{ENV} directly;
  the four \file{SPH\_*} knobs are read the same way inside \file{SimParams}.
\item \textbf{Driver settings} (box, $s_0$, $t_{\rm end}$, seeds, record interval, boost, \dots) are
  read once, at the top of the driver, through \file{setting}, into one named tuple.
\end{itemize}

%=============================================================================
\section{The run record}
\label{sec:record}

Every driver that writes data writes \file{<run>.toml} next to it, generated from the
structs and the settings read, never typed in:
\begin{center}
\begin{tabular}{lp{0.5\textwidth}}
\hline
section & contents \\ \hline
\file{[sim]}, \file{[eos]}, \file{[stochastic]} & every field of the struct \\
\file{[driver]} & every \file{setting} read, with the value used and its default \\
\file{[start]} & how the initial state was made, the seeds \\
\file{[provenance]} & commit (with \file{+dirty} and a hash of the diff), Julia, threads, host, Slurm job and task, date \\
\file{[outcome]} & steps done, complete, wall time, KE/eq, $\Delta H/H$, monitor totals, guards \\
\hline
\end{tabular}
\end{center}
It is written at the start, rewritten at each checkpoint and at the end (\file{[outcome]}
changes, the rest does not).  TOML is in the Julia standard library (read and write) and
Python reads it with \file{tomllib}.  The one-row \file{.csv} stays for
\file{nersc/collect\_runs.jl}.  A unit test checks that the record holds every field of every
settings struct and every \file{SR\_}/\file{SPH\_} name read.  The log (\file{slurm/*.out})
keeps the monitor line every 5\% of the run and the summary.

%=============================================================================
\section{The test suite}
\label{sec:tests}

Three kinds of test, all under \file{test/}, run by one script, \file{test/run.sh}, locally or
on one NERSC debug node (\file{nersc/unittests.sh} and \file{nersc/dtests.sh} become wrappers
of it):
\begin{itemize}
\item \textbf{Unit} (\file{test/*.jl}, as now; seconds each): the kernel, cell list, EOS,
  $\kappa$, change of representation, guards, checkpoint, equilibrium start, correlators.
\item \textbf{Regression} (\file{test/regress/}): \file{testregress\_step.jl} plus a short
  fixed-seed run of each production driver.  The reference is not stored: \file{run.sh
  regress <ref>} builds it from a commit (default \file{main}) in a git worktree and compares
  bit for bit.  This is the check of \Sect{sec:rule}.
\item \textbf{Physics} (\file{test/physics/}; minutes): the checks now in \file{examples/test*.jl}
  and \file{shearwave\_decay.jl} --- sound, viscous sound in 1D and 2D, Gaussian diffusion in 1D
  and 2D, the shear wave --- each with a pass/fail criterion (a measured rate or speed against
  the analytic one, within a stated tolerance set from the present results), run for the
  three schemes.  Figures stay optional (\file{SR\_FIGDIR}).
\end{itemize}
Each test writes one line to \file{results.toml} in the output directory (name, scheme, pass,
the measured numbers, time); \file{run.sh} prints the table and exits nonzero on a failure.

%=============================================================================
\section{Order of the steps}
\label{sec:order}

\begin{enumerate}
\item \file{src/load.jl}, and the scripts use it.
\item The library split by content (\file{ideal.jl}, \file{pointwise.jl}, \file{viscous.jl}), and the new names of the deterministic step.
\item \file{src/legacy/}: the superseded code and its users.
\item The rename to ``stochastic'' (code, drivers, notes).
\item The kernel names.
\item \file{test/run.sh} and the regression from a worktree (needed before the steps that
  touch drivers).
\item \file{StochOptions} and \file{setting}; the drivers read their settings once.
\item The TOML record and its test.
\item The physics tests with their criteria, moved from \file{examples/}.
\end{enumerate}
Steps 1--6 move or rename code only.  Steps 7--9 change how settings reach the code, not the numbers:
the regression check stays bit-identical through all of them.

\end{document}
